\documentclass[12pt,a4paper,reqno, aps, superscriptaddress]{amsart}

% \usepackage[colorlinks=true,citecolor=blue,urlcolor=blue]{hyperref}

\usepackage[T1]{fontenc}
\usepackage[utf8]{inputenc}

\usepackage{graphicx}
\usepackage[textheight=20 cm, textwidth=15cm, headheight=50pt, headsep=10pt, footskip=32pt]{geometry}

\usepackage{amsfonts}
\usepackage{bm}
\usepackage{amssymb}
\usepackage{amsmath}
\usepackage{amsthm}
\usepackage[foot]{amsaddr}
\usepackage{mathtools}
\usepackage{dsfont}

\usepackage{physics}
\usepackage{tikz}
\input{qcircuit.tex}

\usepackage{caption}
\usepackage{subcaption}
\usepackage[normalem]{ulem} 
\usepackage[colorinlistoftodos]{todonotes}
\usepackage[table]{xcolor}
\usepackage{orcidlink}
\usepackage{diagbox}
\usepackage{booktabs}
\usepackage{dcolumn}
\usepackage{appendix}

%%%%%%%%%%%%%%%%%%%%%%%%%%%%%

\newtheorem{theorem}{Theorem}
\newtheorem{proposition}{Proposition}
\newtheorem{remark}{Remark}

\newcommand{\Id}{{\rm 1\hspace{-0.9mm}l}}
\newcommand{\C}{\ensuremath{\mathbb{C}}}
\newcommand{\PP}{\ensuremath{\mathcal{P}}}
\newcommand{\QQ}{\ensuremath{\mathcal{Q}}}
\newcommand{\R}{\ensuremath{\mathbb{R}}}
\renewcommand{\AA}{\ensuremath{\mathcal{A}}}
\newcommand{\BB}{\ensuremath{\mathcal{B}}}
\newcommand{\p}{\ensuremath{p_{\mathrm{I}}}}
\newcommand{\pp}{\ensuremath{p_{\mathrm{II}}}}
\newcommand{\proj}[1]{\ensuremath{\ketbra{#1}{#1}}}

%%%%%%%%%%%%%%%%%%%%%%%%%%%%%

\title[Quantum key distribution]{High secret key rate QKD protocols}

\author{Ryszard Kukulski, Paulina Lewandowska, Norbert Luchowski, Łukasz 
Pawela, Zbigniew Puchała}

%%%%%%%%%%%%%%%%%%%%%%%%%%%%%

\begin{document}
\maketitle

\section{QKD protocol}

Here we present full QKD protocol overview, where each steps will be explained 
in details in the sections below:
\begin{enumerate}
\item Quantum transmission,
\item Post-selection,
\item Security check,
\item Secure key agreement.
\end{enumerate}


\subsection{Quantum transmission (see Fig.~\ref{fig-protocol})}

\begin{figure}[ht!]
	\begin{align*}
	\begin{aligned}
	\Qcircuit @C=1em @R=.7em {
		p_{i,a} \!\!\!\to & \prepareC{\ket{\psi_{i,a}}}	 & 
		\multigate{1}{E} 
		& \gate{\Omega} & \measureD{\phi_{j,b}}[0] \cw   \\&
		  & \pureghost{E} & \qw & \qw & \gate{\Upsilon_{i}} & \measureD{e}[0] 
		  \cw
	}
	\end{aligned}
	\end{align*}
    \caption{The general setup of QKD protocol we consider.}
    \label{fig-protocol}
\end{figure}

Alice and Bob perform a sequence of quantum state transmission experiments. 
Each of them consists of Alice sending a quantum state to Bob, whose performs a 
quantum measurement. 

In details, Alice sends a pure quantum state $\ket{\psi_{i,a}} \in \C^2$ that 
is 
randomly generated with a probability distribution $p_{i,a}$, where 
$i=1,\ldots,N$ and $a=0,1$. Bob, whereas, prepares POVM $\Omega = 
\{\proj{\phi_{j,b}}\}_{j,b}$ for $j=1,\ldots,N$ 
and $b=0,1$. He reads $\ket{\phi_{j,b}}$ as a result of the measurement.

\subsection{Post-selection}

For each experiment Alice and Bob communicate publicly values of $i$ and $j$. 
They accept the result of given experiment if $i=j$. Otherwise, the result is 
discarded. 

\subsection{Security check}
At this point Alice and Bob have a sequence of accepted experiments' results. 
For each of them, Alice associates the state $\ket{\psi_{i,a}}$ with a raw key 
bit 
value $a$. Bob meanwhile associates the measurement result $\ket{\phi_{i,b}}$ 
with a 
raw key bit value $b$.

Alice publicly announces the values of her bit-string from randomly chosen 
subset that is logarithmic in size with respect to the length of the bit-string.
Bob announces the values of his bit-string for the same subset. They estimate 
percentage of bits that values agree. If this percentage is greater than 
acceptable 
ratio $x \in (0.5,1]$, then the protocol is accepted. Otherwise, the whole 
procedure is aborted.

\subsection{Secure key agreement}
Alice and Bob use the remaining parts of their 
bit-strings to extract secure random key that both agree on by applying Error 
Correction and Privacy Amplification methods.


\section{Eve attacks}

For each experiment, the potential Eve attack is described 
by a quantum channel $E$ and conditional quantum channel $\Upsilon_{i}$, which 
is a POVM measurement returning bit $e$ (see 
Fig.~\ref{fig-protocol}).

Eve targets can be:
\begin{enumerate}
    \item Reducing the probability that $i=j$, which will reduce key ratio,
    \item Reducing the percentage of bits that Alice and Bob agree on, which 
    will lead to protocol abortion or will reduce key ratio,
    \item Increasing the knowledge about secret key by using results $e$, which 
    can break the security of the key or reduce key ratio.
\end{enumerate}

\section{Single experiment measures}
\begin{itemize}
    \item The probability of success post-selection $S$: $\PP(S)$.
    \item The probability of raw key bit agreement $\PP(B=A|S)$.
    \item The probability of eavesdropping the raw $\PP(E=A|S)$ (for Alice's 
    bit-string) or $\PP(E=B|S)$ (for Bob's bit-string).
\end{itemize}

\section{Noise model and security constraints}
\subsection{Noise model}
We need to model noisy quantum communication channel, errors that can occur 
during state preparation or measurement and influence of Eve. All these effects 
can be captured by the notion of a single quantum channel $\Phi$. Ideally, 
$\Phi=\Id$ but the more the noise is intensive or Eve attacks are harder the 
more
$\Phi$ drifts away from $\Id$. To measure the intensity of the noise, we 
will 
use quantum channel fidelity between Choi-Jamio{\l}kowski operators of $\Phi$ 
and $\Id$ given by 
\begin{equation}
\frac14 \bra{\Id} J_\Phi \ket{\Id} \in [0,1].
\end{equation}
We say that $\Phi$ intensity is at most $\epsilon \in [0,1]$ (and write $\Phi 
\sim \epsilon$), when $\frac14 
\bra{\Id} J_\Phi \ket{\Id} \ge 1-\epsilon$.

\subsection{Security}
Taking $\epsilon$ as given for QKD purposes we can compute the minimal 
probability that raw key bits will agree. However, in QKD protocol during 
Security check we can estimate this value experimentally. The truth value can 
be lower than the threshold that we set and hence, we abort the protocol or it 
can be within the range. However, to obtain unconditional security of our QKD 
we can only assume that QBER is not greater than the threshold we defined.

\section{Problem definition}
For a given $\epsilon \in [0,1]$ and QKD protocol we can compute secure key 
ratio $R$ which tells how many secure bits Alice and Bob can agree per single 
quantum experiment. We assume that Secure key agreement is done by 
direct-reconciliation (without Advantage Distillation) and hence the ratio for 
EC and PA is given by $\max(I(A;B) - I(A;E),I(A;B) -I(B;E)) = 
\max(H(A|E)-H(A|B),H(B|E)-H(B|A))$. Due to possibility of public key 
symmetrization, we get $H(A|B)=H(B|A)=H_2(\PP(B=A|S))$. On the other hand 
$H(P|E)=H_2(\PP(E=P|S))$, where $P \in \{A,B\}$. Taking into 
account the probability of success $\PP(S)$ we get
\begin{equation}
R = \PP(S)(\max_P H_2(\PP(E=P|S)) - H_2(\PP(B=A|S))).
\end{equation}

Let us put an important assumption for our QKD protocol. If $\epsilon=0$ 
then the Secure key agreement should work at its best, which means parties' 
raw keys are perfectly correlated, $\PP(B=A|S)=1$. Moreover, there is no 
leakage of information to Eve, which is $\PP(E=P|S)=0.5$. Right now, we 
would like to define \textbf{the worst case} SKR $R_\epsilon$ 
within the 
noise model with the intensity $\epsilon$. Note that $\PP(S)$ and $\PP(B=A|S)$ 
can be optimized within the scope of the model. Although, to guarantee 
unconditional security, the value of $\PP(E=P|S)$ will be optimized according 
to the threshold given by $\PP(B=A|S)$.

We define
\begin{itemize}
\item $\PP_\epsilon(S) = \min_{\Phi \sim \epsilon} \PP(S)$.
\item $\PP_\epsilon(B=A|S) = \max\left\{\min_{\Phi \sim \epsilon} 
\PP(B=A|S), \frac12 \right\}$.
\item $\PP_\epsilon(E=P|S) = \max_{E} \{\PP(E=P|S): \PP(B=A|S) \ge 
\PP_\epsilon(B=A|S)\}$, where the maximization is over all Eve strategies.
\end{itemize}

\subsection{Goal}
For a given $\epsilon \in [0,1]$ and the QKD protocol we define the worst case 
SKR as
\begin{equation}
R_\epsilon = \PP_\epsilon(S)(H_2(\min_P \PP_\epsilon(E=P|S)) - 
H_2(\PP_\epsilon(B=A|S))).
\end{equation}

Our goal is to find (numerically / analytically)
\begin{equation}
R^*_\epsilon = \sup_{\text{QKD}}R_\epsilon.
\end{equation}

\section{Results}
\todo[inline]{!!! Łukasz i Norbert}

\subsection{QKD parameters:}
\begin{itemize}
\item $N \in \mathbb{N}$ - the number of pairs,
\item $(p_{i,a})_{i,a}$, where $i=1,\ldots,N$ and $a=0,1$. The vector satisfies 
$p_{i,a} > 0$ for all $i,a$ and $\sum_{i,a} p_{i,a} = 1$,
\item $(\ket{\psi_{i,a}})_{i,a}$, where $i=1,\ldots,N$ and $a=0,1$. The vectors 
$\ket{\psi_{i,a}} \in \C^2$ satisfy $\|\psi_{i,a}\|_2=1$ and pairs of vectors 
$\ket{\psi_{i,0}}, \ket{\psi_{i,1}}$ 
are linearly independent, for all $i$,
\item $(q_i)_i$ for $i=1,\ldots,N$. The vector satisfy $q_i > 0$ for all $i$ 
and \begin{equation}
\sum_{i,a} 
q_i \frac{ p_{i,a \oplus 1}}{p_{i,0}+p_{i,1}} \proj{\psi_{i,a \oplus 1}^\perp} 
\le \Id_2.\end{equation}
\end{itemize}
\subsection{Input:}
$\epsilon \in [0,1]$
\subsection{Output:}
A function that computes $\epsilon \mapsto \sup_{\text{QKD}} R_\epsilon$, where 
$R_\epsilon$ is the minimal SKR for given $\epsilon$ and QKD.

\appendix
\section{Mathematical objects}
The noise channel $\Phi \sim \epsilon$ is defined by its isomorphism $J_\Phi 
\in \mathrm{P}(\AA \otimes \BB)$ 
for which $\frac14 \bra{\Id} J_\Phi \ket{\Id} \ge 1-\epsilon$. Then, the 
probability that Alice sends $\ket{\psi_{i,a}}$ and Bob measures 
$\ket{\phi_{j,b}}$ equals 
$p_{i,a} \tr\left(J_\Phi (\overline{\proj{\psi_{i,a}}} \otimes 
\proj{\phi_{j,b}})\right)$. 
That gives
\begin{equation}
\PP(S) = \sum_{i,a,b} p_{i,a} \tr\left(J_\Phi (\overline{\proj{\psi_{i,a}}} 
\otimes 
\proj{\phi_{i,b}})\right).
\end{equation}
Similarly we have
\begin{equation}
\PP(B=A \wedge S) = \sum_{i,a} p_{i,a} \tr\left(J_\Phi 
(\overline{\proj{\psi_{i,a}}} 
\otimes 
\proj{\phi_{i,a}})\right).
\end{equation}

To compute $\PP_\epsilon(E=P|S)$ we take Eve's actions. We define Eve strategy 
as a set of operators $\{ E_{i,e}\}_{i,e}$ with $i=1,\ldots,N$ and $e=0,1$ with 
network conditions defined as 
\begin{equation}
    E_{i,e} \ge 0, \,\,\, \sum_{e} E_{i,e} = \sum_{e} E_{j,e}, \,\,\, 
    \tr_{2} \left( \sum_{i,e} E_{i,e}   \right) = N \Id_2 .
\end{equation} 
Let us consider interactions of Eve with QKD code.
\begin{equation}
    \PP(S) = \sum_{i,a,b,e} p_{i,a} \tr\left(E_{i,e} 
    (\overline{\proj{\psi_{i,a}}} 
\otimes 
\proj{\phi_{i,b}})\right).
\end{equation}
Similarly we have
\begin{equation}
\PP(B=A \wedge S) = \sum_{i,a,e} p_{i,a} \tr\left(E_{i,e}
(\overline{\proj{\psi_{i,a}}} 
\otimes 
\proj{\phi_{i,a}})\right)
\end{equation}
and
\begin{equation}
\begin{split}
\PP(E=A \wedge S) &= \sum_{i,a,b} p_{i,a} \tr\left(E_{i,a} 
    (\overline{\proj{\psi_{i,a}}} 
\otimes 
\proj{\phi_{i,b}})\right),\\
\PP(E=B \wedge S) &= \sum_{i,a,b} p_{i,a} \tr\left(E_{i,b} 
    (\overline{\proj{\psi_{i,a}}} 
\otimes 
\proj{\phi_{i,b}})\right).
\end{split}
\end{equation}

\section{Optimization problems}
We start by reducing the number of parameters 
which define the given QKD protocol by using the assumption for 
$\epsilon=0$. If Alice 
generated the state $\ket{\psi_{i,a}}$ and Bob measured $\ket{\phi_{i,b}}$ that 
means $a=b$. 
In other words, $\ket{\phi_{i,b}} \propto \ket{\psi_{i,b\oplus 1}^\perp}$. 
Moreover, for each $i$ 
there can not be any bias between $0$ and $1$. That indicate $\ket{\phi_{i,b}} 
= \sqrt{\frac{ p_{i,b \oplus 1}}{p_{i,0}+p_{i,1}}q_i}\ket{\psi_{i,b\oplus 
1}^\perp}$ for some $q_i > 0$.
\subsection{$\PP_\epsilon(S)$:}
Let $Q_S = \sum_{i,a,b}  
p_{i,a} \overline{\proj{\psi_{i,a}}} 
\otimes \proj{\phi_{i,b}}$.
\begin{equation}
\begin{split}
\PP_\epsilon(S) &= \min_{\Phi \sim \epsilon} \PP(S)\\
&=\min_{J_\Phi} \left( \tr\left(J_\Phi Q_S\right): J_\Phi \ge 0, \tr_2(J_\Phi) 
= \Id_2,
\frac14 \bra{\Id} J_\Phi \ket{\Id} \ge 1-\epsilon \right).
\end{split}
\end{equation}

\subsection{$\PP_\epsilon(B=A|S)$:}
Let $W_S = \sum_{i,a}  
p_{i,a} \overline{\proj{\psi_{i,a}}} 
\otimes \proj{\phi_{i,a}}$.
\begin{equation}
\begin{split}
&\PP_\epsilon(B=A|S)\\ =& \min_{\Phi \sim \epsilon} \PP(B=A|S)\\
=&\min_{J_\Phi} \left( \frac{\tr\left(J_\Phi W_S\right)}{\tr\left(J_\Phi 
Q_S\right)}: J_\Phi \ge 0, \tr_2(J_\Phi) 
= \Id_2,
\frac14 \bra{\Id} J_\Phi \ket{\Id} \ge 1-\epsilon \right)\\
=&\min_{J_\Phi} \left( \frac{\tr\left(J_\Phi W_S\right)}{\tr\left(J_\Phi 
Q_S\right)}: J_\Phi \ge 0, \tr_2(J_\Phi) 
\propto \Id_2,
\bra{\Id} J_\Phi \ket{\Id} \ge \tr(J_\Phi)(2-2\epsilon) \right)\\
=&\min_{J_\Phi} \left( \tr\left(J_\Phi W_S\right): J_\Phi \ge 0, 
\tr_2(J_\Phi) 
\propto \Id_2,
\bra{\Id} J_\Phi \ket{\Id} \ge \tr(J_\Phi)(2-2\epsilon), \tr\left(J_\Phi 
Q_S\right)=1 \right).
\end{split}
\end{equation}



\subsection{$\PP_\epsilon(E=P|S)$:} Let $x=\PP_\epsilon(B=A|S)$ be the 
theoretical threshold value. Note that $\PP(S)$ and $\PP(E=P \wedge S)$ are 
linear functions of $E$.
\begin{equation}
\begin{split}
&\PP_\epsilon(E=P|S)\\ =& \max_{E} \{\PP(E=P|S): \PP(B=A|S) \ge 
x\}\\
=& \max_{E} \{\frac{\PP(E=P \wedge S)}{\PP(S)}: \frac{\PP(B=A \wedge 
S)}{\PP(S)} \ge x, 
\begin{cases} E_{i,e} \ge 0,\\ \sum_{e} E_{i,e} = \sum_{e} E_{j,e},
\\    \tr_{2} \left( \sum_{i,e} E_{i,e} \right) = N \Id_2 \end{cases} \}\\
=& \max_{E} \{\frac{\PP(E=P \wedge S)}{\PP(S)}: \frac{\PP(B=A \wedge 
S)}{\PP(S)} \ge x, 
\begin{cases} E_{i,e} \ge 0,\\ \sum_{e} E_{i,e} = \sum_{e} E_{j,e},
\\    \tr_{2} \left( \sum_{i,e} E_{i,e} \right) \propto N \Id_2 \end{cases} \}\\
=& \max_{E} \{\PP(E=P \wedge S): \PP(B=A \wedge 
S) \ge x, \PP(S)=1,
\begin{cases} E_{i,e} \ge 0,\\ \sum_{e} E_{i,e} = \sum_{e} E_{j,e},
\\    \tr_{2} \left( \sum_{i,e} E_{i,e} \right) \propto N \Id_2 \end{cases} \}\\
\end{split}
\end{equation}

\subsection{$R^*_\epsilon$:} ?



\bibliographystyle{unsrt} 
\bibliography{qkd.bib}
\end{document}
